\subsection{Propiedades del paso al límite proyectivo.}
\label{subsection:IV.19.8}
\label{IV.19.8}

En este número, las notaciones y convenciones sobre los límites proyectivos son las de (8.5.1) y (8.8.1).

\begin{proposition}[19.8.1]
\label{IV.19.8.1}
Supongamos que los morfismos de transición $S_\mu\to S_\lambda$ $(\lambda\leq\mu)$ sean planos y, además, que se cumpla una de las dos hipótesis siguientes:
\begin{enumerate}
  \item[$1^\circ$] Los preesquemas $S_\lambda$ son localmente noetherianos.
  \item[$2^\circ$] Los morfismos de transición $S_\mu\to S_\lambda$ son sobreyectivos (y, por tanto, fielmente planos).
\end{enumerate}

\oldpage[IV-4]{220}

\emph{Bajo estas condiciones:}
\begin{enumerate}
  \item[(i)] Supongamos que $S_\alpha$ sea cuasicompacto. Sea $\sh F_\alpha$ un $\sh O_{S_\alpha}$-módulo cuasicoherente, que suponemos además de tipo finito cuando se cumple la hipótesis $1^\circ$. Sea $(f_{i\alpha})_{1\leq i\leq n}$ una sucesión finita de secciones de $\sh F_\alpha$ sobre $S_\alpha$. Para que la sucesión de secciones $f_i$ del $\sh O_S$-módulo $\sh F$ sobre $S$, correspondiente a las $f_{i\alpha}$ $(1\leq i\leq n)$, sea $\sh F$-regular, es necesario y suficiente que exista $\lambda\geq\alpha$ tal que la sucesión de secciones $f_{i\lambda}$ de $\sh F_\lambda$ sobre $S_\lambda$ sea $\sh F_\lambda$-regular.
  \item[(ii)] Supongamos que $X_\alpha$ sea cuasicompacto, y sea $j_\alpha:X_\alpha\to S_\alpha$ una inmersión, que suponemos localmente de presentación finita cuando se cumple la hipótesis $2^\circ$. Entonces, para que la inmersión correspondiente $j:X\to S$ sea regular, es necesario y suficiente que exista $\lambda\geq\alpha$ tal que $j_\lambda:X_\lambda\to S_\lambda$ sea regular.
\end{enumerate}
\end{proposition}

\begin{enumerate}
  \item[(i)] Si $p_\lambda:S\to S_\lambda$ es la proyección canónica, se tiene
  \[
    \sh F/\Bigl(\sum_{j=1}^{i-1}f_j\sh F\Bigr)
    =p_\lambda^*\Bigl(\sh F_\lambda/\Bigl(\sum_{j=1}^{i-1}f_{j\lambda}\sh F_\lambda\Bigr)\Bigr),
  \]
  por tanto, nos reducimos al caso $n=1$, en cuyo caso se suprime el índice $i$. La suficiencia de la condición resulta de que $p_\lambda$ es plano (8.3.8) y de ($0$, 15.1.5). En el caso $2^\circ$, $p_\lambda$ es fielmente plano (8.3.8), y la necesidad de la condición resulta también de ($0$, 15.2.5), con $\lambda=\alpha$. En el caso $1^\circ$, denotemos por $\sh N_\lambda$ (resp. $\sh N$) el núcleo del endomorfismo $\sh F_\lambda\to\sh F_\lambda$ (resp. $\sh F\to\sh F$) dado por la multiplicación por $f_\lambda$ (resp. por $f$); como $\sh F_\lambda$ es coherente por hipótesis, también lo es $\sh N_\lambda$, de modo que la hipótesis $\sh N=0$ implica $\sh N_\lambda=0$ para algún $\lambda\geq\alpha$ en virtud de (8.5.8, (ii)).
  \item[(ii)] Como $p_\lambda$ es plano, la suficiencia de la condición resulta de (19.1.5, (ii)). Para demostrar su necesidad, observemos que, como $j_\alpha(X_\alpha)$ es cuasicompacto, está contenido en un abierto cuasicompacto de $S_\alpha$, y podemos reducirnos al caso en que $S_\alpha$ también sea cuasicompacto y $j_\alpha$ una inmersión cerrada; así, la imagen de $X_\alpha$ está definida por un ideal cuasicoherente $\sh J_\alpha$ de $\sh O_{S_\alpha}$, que además podemos suponer de tipo finito en los casos $1^\circ$ y $2^\circ$ ($j_\alpha$ es localmente de presentación finita en el caso $2^\circ$); la imagen de $X_\lambda$ (resp. $X$) en $S_\lambda$ (resp. $S$) está entonces definida por $\sh J_\lambda=\sh J_\alpha\sh O_{S_\lambda}$ (resp. $\sh J=\sh J_\alpha\sh O_S$), que sigue siendo de tipo finito. Además, teniendo en cuenta (8.2.11), podemos suponer que $\sh J$ está engendrado por una sucesión regular $(f_i)_{1\leq i\leq n}$ de secciones sobre $S$, que por tanto definen un homomorfismo sobreyectivo $u:\sh O_S^n\to\sh J$. Teniendo en cuenta (8.5.2, (i)) y (8.5.7), existen $\lambda\geq\alpha$ y un homomorfismo sobreyectivo $u_\lambda:\sh O_{S_\lambda}^n\to\sh J_\lambda$ tales que $u=p_\lambda^*(u_\lambda)$; por tanto, los $f_i$ son las imágenes canónicas de secciones $f_{i\lambda}$ de $\sh J_\lambda$ sobre $S_\lambda$ que engendran este ideal. En virtud de (i), existe entonces $\mu\geq\lambda$ para el cual la sucesión $(f_{i\mu})$ es $\sh O_{S_\mu}$-regular, y por tanto la inmersión $j_\mu$ es regular.
\end{enumerate}

\begin{proposition}[19.8.2]
\label{IV.19.8.2}
Sea $X_\alpha$ un preesquema $S_\alpha$ cuasicompacto y localmente de presentación finita.
\begin{enumerate}
  \item[(i)] Sean $\sh F_\alpha$ un $\sh O_{X_\alpha}$-módulo cuasicoherente de presentación finita y $(f_{i\alpha})_{1\leq i\leq n}$ una sucesión de secciones de $\sh O_{X_\alpha}$ sobre $X_\alpha$. Para que la sucesión de secciones $f_i$ de $\sh O_X$ sobre $X$, correspondiente a las $f_{i\alpha}$, sea $\sh F$-transversalmente regular relativamente a $S$ (19.2.1), es necesario y suficiente que exista $\lambda\geq\alpha$ tal que la sucesión de secciones $f_{i\lambda}$ de $\sh O_{X_\lambda}$ sobre $X_\lambda$ sea $\sh F_\lambda$-transversalmente regular relativamente a $S_\lambda$.
  \item[(ii)] Sea $Y_\alpha$ un preesquema $S_\alpha$ cuasicompacto, y sea $j_\alpha:Y_\alpha\to X_\alpha$ una inmersión $S_\alpha$ localmente

\oldpage[IV-4]{221}

  de presentación finita. Para que la inmersión $S$ correspondiente $j:Y\to X$ sea transversalmente regular relativamente a $S$, es necesario y suficiente que exista $\lambda\geq\alpha$ tal que $j_\lambda:Y_\lambda\to X_\lambda$ sea transversalmente regular relativamente a $S_\lambda$.
\end{enumerate}
\end{proposition}

\begin{enumerate}
  \item[(i)] La suficiencia de la condición resulta de ($0$, 15.1.15). Para demostrar que es necesaria, bastará mostrar que todo punto $x\in X$ posee un entorno abierto $V(x)$ para el cual existen un índice $\lambda=\lambda(x)\geq\alpha$ y un entorno abierto $V_\lambda$ de la imagen $x_\lambda$ de $x$ en $X_\lambda$, tales que $V(x)$ sea la imagen inversa de $V_\lambda$ y la sucesión de las restricciones a $V_\lambda$ de los $f_{i\lambda}$ sea $(\sh F_\lambda|V_\lambda)$-transversalmente regular relativamente a $S_\lambda$; en efecto, bastará entonces aplicar (8.3.4) al abierto $U_\lambda$ de $X_\lambda$ definido como el mayor abierto en el que las restricciones a $U_\lambda$ de los $f_{i\lambda}$ forman una sucesión $(\sh F_\lambda|U_\lambda)$-transversalmente regular relativamente a $S_\lambda$.

  Sea $s$ la imagen de $x$ en $S$, y para todo $\lambda\geq\alpha$ sea $s_\lambda$ la imagen de $s$ en $S_\lambda$. Entonces la fibra $X_s$ es el límite proyectivo de las fibras $(X_\lambda)_{s_\lambda}$ y es localmente noetheriana, puesto que el morfismo $X\to S$ es localmente de presentación finita; como además $(X_\alpha)_{s_\alpha}$ es cuasicompacto y los morfismos de transición $(X_\mu)_{s_\mu}\to(X_\lambda)_{s_\lambda}$ son planos (pues también lo es $\Spec(k(s_\mu))\to\Spec(k(s_\lambda))$), podemos aplicar a las secciones $f_i\otimes1$ de $\sh F_s=\sh F\otimes_{\sh O_S}k(s)$ sobre $X_s$ el resultado de (19.8.1, (i)); de aquí se obtiene un $\lambda\geq\alpha$ tal que las secciones $f_{i\lambda}\otimes1$ de $(\sh F_\lambda)_{s_\lambda}=\sh F_\lambda\otimes_{\sh O_{S_\lambda}}k(s_\lambda)$ sobre $(X_\lambda)_{s_\lambda}$ forman una sucesión $(\sh F_\lambda)_{s_\lambda}$-regular. Además, por (11.2.6), podemos suponer que $\sh F_\lambda$ es $S_\lambda$-plano en el punto $x_\lambda$. Se deduce entonces de (11.3.8) que existe un entorno abierto $V_\lambda$ de $x_\lambda$ en $X_\lambda$ que satisface las condiciones requeridas.

  \item[(ii)] La suficiencia de la condición resulta de (19.2.7, (ii)). Para demostrar su necesidad, basta todavía mostrar que todo punto $x\in j(Y)$ admite un entorno abierto $V(x)$ para el cual existen un índice $\lambda=\lambda(x)\geq\alpha$ y un entorno abierto $V_\lambda$ de la imagen $x_\lambda$ de $x$ en $X_\lambda$, tales que $V(x)$ sea la imagen inversa de $V_\lambda$ y la restricción $j_\lambda^{-1}(V_\lambda)\to V_\lambda$ de $j_\lambda$ sea transversalmente regular relativamente a $S_\lambda$. Sea $s$ la imagen de $x$ en $S$, y para todo $\lambda\geq\alpha$ sea $s_\lambda$ la imagen de $s$ en $S_\lambda$. Como $X_s$ (resp. $Y_s$) es el límite proyectivo de las fibras $(X_\lambda)_{s_\lambda}$ (resp. $(Y_\lambda)_{s_\lambda}$), podemos utilizar, como en (i), (19.8.1, (ii)); y como, por hipótesis, la inmersión $Y_s\to X_s$ es regular, existe un índice $\lambda(s)$ tal que, para $\lambda\geq\lambda(s)$, la inmersión $(Y_\lambda)_{s_\lambda}\to(X_\lambda)_{s_\lambda}$ es regular. Además, en virtud de (19.2.4), existe un entorno abierto cuasicompacto $W(x)$ de $x$ en $X$ tal que los morfismos estructurales $W(x)\to S$ y $Y\cap W(x)\to S$ son planos; aplicando (11.2.6) y (8.2.11), podemos suponer que existen $\lambda\geq\lambda(s)$ y un entorno abierto cuasicompacto $W(x_\lambda)$ de $x_\lambda$ en $X_\lambda$, tales que $W(x)$ sea la imagen inversa de $W(x_\lambda)$ y las restricciones a $W(x_\lambda)$ y a $Y_\lambda\cap W(x_\lambda)$, respectivamente, de los morfismos estructurales $X_\lambda\to S_\lambda$ y $Y_\lambda\to S_\lambda$ sean planas. Concluimos, pues, de (19.2.4) que existe un entorno abierto $V_\lambda$ de $x_\lambda$ que responde a la cuestión.
\end{enumerate}

\begin{remark}[19.8.3]
\label{IV.19.8.3}
Las demostraciones precedentes muestran que los enunciados de (19.8.1) y (19.8.2) siguen siendo válidos cuando se sustituyen las condiciones relativas a la regularidad o a la regularidad transversal en todo $X$ (resp. $X_\lambda$) por esas condiciones en un entorno de un punto $x\in X$ (resp. en un entorno de $x_\lambda$, proyección de $x$).
\end{remark}
